\section{Full exact-ReLU algebra}
\label{app:relu}

This appendix expands the proof obligations behind
\cref{lem:relu-complementarity,lem:relu-polynomialization}.

\subsection{Case partition}

The equation \(a(a-z)=0\) gives the algebraic alternatives
\(a=0\) and \(a=z\).  The two weak inequalities select the valid half of each
branch:
\begin{center}
\begin{tabular}{cccc}
\toprule
Algebraic branch & Consequence of inequalities & ReLU regime & Included boundary\\
\midrule
\(a=0\) & \(z\le0\) & inactive & \(z=0\)\\
\(a=z\) & \(z\ge0\) & active & \(z=0\)\\
\bottomrule
\end{tabular}
\end{center}
At \(z=a=0\), both algebraic branches meet, but they represent the same ReLU
assignment.  The complementarity description therefore has no spurious
\((z,a)\)-point.

\subsection{Square-slack fibers}

If \(a>0\), the equation \(a=u^2\) has two real slack witnesses.  If \(a=0\),
it has the unique witness \(u=0\).  The same holds for \(a-z=v^2\).  Hence the
projection of the slack variety to \((z,a)\) is exactly the ReLU graph, although
the map is generally many-to-one.

For one gate, the equality-only encoding adds two variables and three
equalities, all of degree at most two.  For \(G\) ReLU gates and \(N\) samples,
the construction adds \(2GN\) real variables and \(3GN\) polynomial
equalities.  With an explicit architecture and dataset, this is polynomial in
the input length.

\subsection{Boundary Jacobian}

For
\begin{align}
 h_1&=a-u^2,\\
 h_2&=a-z-v^2,\\
 h_3&=a(a-z),
\end{align}
the Jacobian in \((z,a,u,v)\) at the boundary origin is
\begin{equation}
  \begin{bmatrix}
   0&1&0&0\\
   -1&1&0&0\\
   0&0&0&0
  \end{bmatrix}.
\end{equation}
The complementarity row vanishes and the slack derivatives vanish.  Thus
existential square polynomialization provides exact semantics but does not
remove singular behavior from the equality presentation.

\subsection{Formal verification correspondence}

The Lean theorem
\path{exactReLU_iff_complementarity} proves
\cref{lem:relu-complementarity}.  The theorem
\path{exactReLU_iff_polynomialized} proves
\cref{lem:relu-polynomialization}.  Separate lemmas cover active, inactive, and
boundary regimes.  These results concern the scalar gate and are composed
sample-by-sample in the mathematical encoding.
